\section{The Seers' Later Lives and the Reliability Problem}

\subsection{Why this section exists}

The 1851 act judged an apparition, not two biographies. But La
Salette differs from the other great nineteenth-century French
apparition in one respect that cannot be managed by tact: at Lourdes
the seer entered religious life, kept silence, and died young; at La
Salette both seers lived long, unsettled public lives, one of them
published an ever-growing revelation, and both were the object of
polarized apologetic and hostile writing within their own lifetimes.
A study that suppressed this would misrepresent the dossier. A study
that made it a scandal would be doing the hostile literature's work.
The discipline adopted here is the one the sources permit: report
what documented witnesses record, attribute contested explanations
to those who offer them, and keep the categories separate. Nothing
in the later lives touches the object of the 1851 judgment, and
nothing in the 1851 judgment protects the later lives from
assessment.

\subsection{Maximin Giraud (1835--1875)}

The Sanctuary's own portrait of Maximin is affectionate and
unflattering in the same breath: during the apparition he amused
himself spinning his hat on his stick and pushing pebbles toward the
Lady's feet, ``and so he remained all his life''---answering
investigators simply but tartly, cordial when he felt loved,
mischievous when someone tried to appropriate him, improvident in
business. Its summary of his career is a list of failures to settle:
the school at Corps, the minor seminary of the Rondeau, a country
presbytery, the Grande Chartreuse, the seminary of Aire-sur-Adour,
the hospice of Le Vésinet, the college of Tonnerre, a household near
Versailles where there was talk of sending him to study medicine,
enlistment as a pontifical zouave and return to Paris six months
later, then a failed liqueur partnership that left him again without
resources.\endnote{Sanctuary, ``Les témoins: Mélanie et Maximin''
(archived capture, \S2 note~3). The Sanctuary's phrase for his
upbringing---``Maximin va pousser comme une herbe sauvage''---and its
judgment ``il demeure alors à Corps, qu'il aurait mieux fait de ne
jamais quitter'' are the custodian's own.}

Two documented acts of his adult life matter for the dossier. He
published a pamphlet, \work{Ma profession de foi sur l'Apparition de
Notre-Dame de La Salette}, in answer to a hostile press article; and
he made a will restating his testimony about the apparition and
leaving his heart to the shrine. He died at Corps on 1 March 1875,
after Communion and a little of the La Salette water, having gone up
the mountain a last time and told the story on the
spot.\endnote{Sanctuary, ``Les témoins.'' The pamphlet is identified
by title from that page; it was not consulted, and nothing about its
contents is asserted here beyond the fact that it was a public
restatement of his testimony (Appendix~A).}

\subsection{The episode at Ars, 1850}

The single most-used argument against Maximin's veracity is his 1850
visit to Jean-Marie Vianney, the curé of Ars, after which it was
said that he had confessed to having lied about the apparition. The
episode is genuinely contested, and this study reports it as
contested. The fullest account it could reach is in the 1854 witness,
which is frankly pro-La Salette and prints Maximin's own answers to
a friend's questions some time later.

On that account: three men of a political circle drew the boy to
Lyons and, under pretence of helping him settle his vocation, took
him to Ars, hoping to obtain his secret. The curé, nearly deaf and
missing most of his teeth, asked whether he had seen the Blessed
Virgin. Maximin says he answered: ``I do not know whether it was the
Blessed Virgin; I saw something---a Lady. But if you know yourself
that it was the Blessed Virgin, you should tell it to all the
people, that they may believe in La Salette.'' Asked directly
whether he had accused himself to the curé of having told lies, he
said he had said that he had sometimes told lies; told he must
retract, he answered that it was too old and not worth the trouble;
and when pressed on what he had meant, he answered ``with a tone of
assurance'': ``I meant my little stories to the Curé of Corps, when I
did not want to say where I was going, or when I did not want to
study my lessons.'' Asked whether the curé of Ars had understood
those stories as referring to the apparition, he answered: ``Why,
yes! he understood it in that way, as that was put in the
newspapers.'' He adds that he was not at confession---``I was at the
confessional, but I had not said my \term{Confiteor}, and I did not
go there to make my confession''---and that, having escaped his
handlers at a Lyons hotel when a priest he knew walked in, he wrote
to the curé of Ars ``to assure him that I had never retracted,
because it is true.''\endnote{Ullathorne 1854, pp.~136--140, printing
the exchange as recorded by ``Mlle Brulais.'' The witness's own
framing---that Maximin went ``in a thoughtless humour'' and came back
in one, and that one of the three men later said he left Ars joyous
and contented---is the witness's, and is reported as such. This
study establishes only that the episode occurred, that Maximin's
reported explanation is on record from 1854, and that the contrary
report was in the newspapers of the day. It does not adjudicate what
was said in the confessional at Ars, and no testimony of Vianney was
reached (Appendix~A).}

The evidential value of the episode is therefore limited in both
directions. It cannot be used to show that the seer recanted,
because the only reasonably early record of what he said is his own
denial together with an innocent explanation. It cannot be used to
show that nothing happened, because the report reached the press and
required an answer. What it does show plainly is a fifteen-year-old
of unstable judgment being handled by adults with an agenda---which
is itself part of the reliability picture, and which the same witness
records that the bishop ended by putting him back in the minor
seminary.

\subsection{Mélanie Calvat (1831--1904)}

The Sanctuary's own account is again the frankest witness available,
and it is worth reading closely because a custodian writing about
its own seer had every incentive to write otherwise. She was
``timid, taciturn, withdrawn,'' yet never hesitated to answer about
the apparition; she spent four years with the Sisters of Providence;
she had little memory and even less aptitude for study than Maximin.
Then the Sanctuary states the problem in its own words: beset by
curiosity, indiscretion, and the pressure of visitors ``avides de
révélations politico-religieuses,'' she ``resists badly the
temptation to play the oracle, taking up again the popular
pseudo-prophecies about the end of times that reappear periodically
in the history of the Church.''

Her itinerary follows: the Carmel of Darlington in England; the
Compassion at Marseille; seventeen years at Castellammare near
Naples, ``writing secrets and a rule for a hypothetical religious
foundation''; a stay at Cannes; Chalon-sur-Saône, where she came into
conflict with the Bishop of Autun for the same reasons; a return to
Italy near Lecce, then Messina; a return to France, in the Allier,
where she finished a mystical autobiography the Sanctuary calls ``de
mauvais aloi.'' On 18--19 September 1902 she passed through La
Salette and told the story of the apparition there. She returned to
southern Italy and died at Altamura, in the province of Bari, the
Sanctuary's page giving 14 December 1904.\endnote{Sanctuary, ``Les
témoins.'' The same page states, of the Castellammare years, that
``le Vatican prie l'évêque du lieu de lui interdire ce genre de
publication mais elle cherche d'autres appuis''---a custodial
statement of a Roman intervention for which this project located no
document, and which is reported here as the Sanctuary's statement
only. The death date is the Sanctuary's; no civil or parish record
was consulted and this study does not adjudicate the datings that
circulate elsewhere (Appendix~A).}

And then the Sanctuary's own conclusion, which is also this study's:
``Pauvre, croyante, pieuse, mais attachée à son propre sens, il est
un point sur lequel Mélanie n'a jamais varié: ce qu'elle avait dit,
comme Maximin, au soir du 19 septembre 1846, dans la cuisine des
Pra, aux Ablandins.''\endnote{Sanctuary, ``Les témoins.'' Editorial
rendering: poor, believing, pious, but attached to her own
judgment---there is one point on which Mélanie never varied: what she
had said, like Maximin, on the evening of 19 September 1846, in the
Pras' kitchen at Les Ablandins. The formulation is exactly the
distinction this study needs and did not have to invent: constancy
about the 1846 recital, and no constancy claimed for anything else.}

\subsection{What biography can and cannot establish}

\begin{boundarytable}{What the biography establishes}{What it does
not establish}
Constancy of both seers about the 1846 recital, from the evening of
the event through interrogations, hostile pressure, offers of money,
and---in Mélanie's case---fifty-six further years, is documented and
is evidence for the recital. & It is not evidence for anything either
of them said, wrote, or published later; the constancy attaches to
the recital, not to the person. \\
Instability, credulity, self-importance, or obstinacy in an adult
life is a real datum about that adult and bears on the weight of
that adult's later claims. & It does not retroactively impeach a
child's testimony given before those traits had any occasion to
operate, and it was not the object of the 1851 judgment. \\
Ecclesiastical conflict with named bishops over publication is
documented and juridically significant. & A conflict about
publication is not a judgment about the apparition, and none of the
bishops involved purported to reopen the 1851 act. \\
\end{boundarytable}

\subsection{Ginoulhiac and the closing of the seers' role}

Jacques-Marie-Achille Ginoulhiac, who succeeded de Bruillard at
Grenoble, spent the 1850s doing two things at once: defending the
1851 judgment against a determined literary opposition, and closing
the door on the seers as continuing organs of revelation.

The defensive acts are dated and specific. The author of a book
against La Salette was cited before the diocesan court and declared
suspended from all ecclesiastical functions on 28 September 1854. A
mandement of 30 September 1854 condemned that book and a related
pamphlet, forbidding all the faithful under pain of excommunication
and all priests under pain of \term{ipso facto} suspension to read,
keep, or circulate them. After consulting the Holy See, a further
mandement of 4 November 1854 condemned a third work, forbidding all
ecclesiastics under pain of suspension to read, keep, or circulate
it. A circular letter to the clergy followed on 19 September 1857
occasioned by later publications on ``the fact and the devotion of
Notre-Dame de la Salette.'' In the same years the controversy also
ran in the civil courts: an 1855 hostile volume records that a
lawsuit was brought at Grenoble by a woman whom two of the
opposition writers had identified in print as the ``white lady'' the
children saw.\endnote{The 1854 acts and the 1857 circular:
\work{Vie de M.~Rousselot, professeur de théologie au grand
séminaire de Grenoble} (Grenoble, 1866), pp.~145--146, read
2026-07-25 in the Internet Archive scan (References)---a Grenoble
diocesan biography, partisan for La Salette and used here for the
dates, titles, and forms of the episcopal acts it reports. The civil
suit: \work{Le miracle de La Salette} (1855), pp.~86--87, an openly
hostile volume, read 2026-07-25 in the Internet Archive scan; this
study records the existence of the litigation on that witness's
authority, establishes nothing about its outcome, was unable to
reach the judgment, and expresses no view on the identification
hypothesis, which the 1851 act's third consideration claims to have
known and weighed (\S4.4).}

The constructive act is a sentence. On 19 September 1855, at the
anniversary, Ginoulhiac stated the principle that has governed La
Salette ever since: «~La mission des bergers est finie, celle de
l'Église commence~»---the shepherds' mission is finished; the
Church's begins.\endnote{Quoted, with that date and attribution, by
the Sanctuary of Notre-Dame de La Salette,
\url{https://lasalette.cef.fr/lhistoire/}, read 2026-07-25, and on
its archived seers page. The Sanctuary presents it as a quotation
from a new inquiry confirming his predecessor's decision. This study
reports it at the custodian's ceiling: the sentence is received in
this polished form, and no printed 1855 text of the address was
reached, so it is cited as the Sanctuary's received formula rather
than as a verified verbatim.}

Whatever its exact original wording, the principle is the correct
juridical description of what happened, and it explains everything in
\S5 without recourse to anyone's bad faith. After 1851 the competent
subject in the La Salette dossier is the Church, acting through
bishops and, when it intervened, through the Holy Office. The seers'
role as witnesses to the 1846 event was complete. Claims that
continued to issue from a seer thereafter belonged to a different
order, and were treated by the competent authorities as such.
